% Abstract
\newpage
\thispagestyle{plain}
\begin{center}
{\Large \textbf{ABSTRACT}}
\end{center}
\vspace{1cm}

Electronic voting systems face core security challenges in simultaneously providing ballot privacy, vote integrity, voter anonymity, coercion resistance, and verifiability. Existing blockchain-based e-voting solutions address only subsets of these requirements, leaving critical gaps in coercion resistance and post-quantum security. This thesis presents a blockchain-based e-voting framework that integrates seven cryptographic protocols into a unified voting pipeline: Paillier homomorphic encryption for encrypted tallying, Pedersen commitments with zero-knowledge $\Sigma$-proofs for vote validity verification, linkable ring signatures for anonymous ballot submission, Self-Masking Deniable Credentials (SMDC) for coercion resistance, Samplable Anonymous Aggregation (SA\textsuperscript{2}) for two-server vote privacy, Kyber768 for post-quantum hybrid encryption, and Damg\aa rd-Jurik-Nielsen threshold decryption for distributed key management. The system is implemented in Go and deployed on Hyperledger Fabric~v2.5~\cite{androulaki2018}. Performance evaluation confirms an $O(n \cdot m)$ total tally complexity ($O(n)$ per candidate) with a constant per-vote per-candidate cost of approximately 6.1~$\mu$s, verified empirically from 100 to 100,000 voters, a complete pipeline time of 70.1~ms per vote, and blockchain throughput of approximately 200~TPS with 100\% success rate at 10,000 transactions. Compared with prior work, the proposed framework reduces the asymptotic tally cost from $O(n \cdot m^2)$ (ISE-Voting) to $O(n \cdot m)$ --- a factor-$m$ improvement --- replaces BP-Vot's $\geq 98\%$ differential-privacy approximation with 100\% exact integer tallies, and adds post-quantum protection at only 5.35~ms per vote ($\approx 7.6\%$ of the pipeline) via the Kyber768 hybrid encryption layer. The novel SMDC mechanism provides computationally indistinguishable real and fake credentials with $k = 5$ slots, enabling coercion resistance in practical blockchain voting contexts. Comparative analysis shows that the proposed framework is the first system to integrate all seven security properties with empirical performance validation on production blockchain infrastructure.

\vspace{1cm}

\noindent
\textbf{Keywords:} blockchain, e-voting, homomorphic encryption, zero-knowledge proofs, ring signatures, coercion resistance, post-quantum cryptography, Hyperledger Fabric, deniable credentials, anonymous aggregation
